% © 2026 Ing. David Jaroš — CC BY-NC-ND 4.0
% Status: GAP-10T-COMPOSITE-FLAT closed [L1], scoped to the gradient-composite scheme.
\documentclass[11pt,a4paper]{article}
\usepackage{amsmath,amssymb,amsthm}
\usepackage[most]{tcolorbox}
\usepackage[margin=2.5cm]{geometry}
\newtheorem{theorem}{Theorem}
\newtheorem{corollary}{Corollary}
\newcommand{\B}{\mathbb B}
\title{Composite Flat Admissibility:\\
The Gradient-Composite Scheme Admits the Affine Representer\\
for All Coefficients}
\author{Ing. David Jaro\v{s}}
\date{26 July 2026}
\begin{document}\maketitle

\begin{abstract}
The effective Palatini variation of the minimal Hilbert--Palatini $+$ kinetic
branch excludes the flat affine representer (flat affine no-go,
\texttt{GAP-10T-FLAT-NOGO}).  We prove the complementary positive statement:
in the \emph{gradient-composite} torsion-free scheme, where the tetrad is the
slice coordinate map of $\partial_\mu\Theta/\sqrt{\mathcal N_0}$ and only
$\Theta$ is varied, the affine representer $\Theta_{\rm aff}$ is a stationary
point of the full action for \emph{all} values of $\Lambda$, $\kappa$, and
$\mathcal N_0$.  The two variational schemes are therefore dynamically
inequivalent at the flat point, and the Palatini no-go does not transfer.
This closes \texttt{GAP-10T-COMPOSITE-FLAT} [L1] and sharpens
\texttt{GAP-10T-DYN}: the surviving minimal continuation is composite.
\end{abstract}

\section{Scheme and statement}
Work on Lorentz-real fields: $\partial_\mu\Theta\in W_L$, tetrad
$e^a{}_\mu$ defined by the slice coordinates of
$\partial_\mu\Theta/\sqrt{\mathcal N_0}$, metric by the central Jordan
identity, spin connection $\mathring\omega(e)$ Levi--Civita (torsion-free
branch), potential $V\equiv0$.  The action is
\begin{equation}
  S[\Theta]
  =\frac1{4\kappa}\!\int\!\epsilon_{abcd}\,e^a\wedge e^b\wedge
    R^{cd}(\mathring\omega(e))
  -\frac{\Lambda}{2\kappa}\!\int\!\sqrt{-g}
  +\frac12\!\int\!\sqrt{-g}\,g^{\mu\nu}
    \langle\partial_\mu\Theta,\partial_\nu\Theta\rangle_\sharp .
  \label{eq:Scomp}
\end{equation}
Only $\Theta$ varies; $e$, $g$, $\sqrt{-g}$, and $\mathring\omega$ vary by the
chain rule.  This is the well-defined gradient-composite scheme; the
self-consistent $D$-composite scheme (connection inside the defining jet)
remains open and is \emph{not} treated here.

\begin{theorem}[Composite flat admissibility]\label{thm:cfa}
$\Theta_{\rm aff}=\Theta_0+\sqrt{\mathcal N_0}\,\bar E_\mu x^\mu$ is a
stationary point of \eqref{eq:Scomp} under compactly supported Lorentz-real
variations, for all $(\Lambda,\kappa,\mathcal N_0)$.
\end{theorem}

\begin{proof}
The Lagrangian density of \eqref{eq:Scomp} is a smooth local function of the
$2$-jet of $\Theta$ in a neighbourhood of the affine jet (the affine tetrad
is nondegenerate, so $\mathring\omega(e)$ is smooth there), and it contains
no undifferentiated $\Theta$ ($V\equiv0$; the scheme replaces $D_\mu$ by
$\partial_\mu$ in the defining jet).  The Euler--Lagrange derivative is
\[
  \mathcal E
  =\frac{\partial L}{\partial\Theta}
  -\partial_\mu\frac{\partial L}{\partial(\partial_\mu\Theta)}
  +\partial_\mu\partial_\nu
   \frac{\partial L}{\partial(\partial_\mu\partial_\nu\Theta)} .
\]
At the affine configuration the $2$-jet is constant in $x$, so both momenta
are constant matrices and their derivatives vanish; the first term vanishes
because $L$ has no undifferentiated $\Theta$.  Hence $\mathcal E=0$.

Concretely, the three pieces vanish for transparent reasons: the kinetic and
$\Lambda$ terms are volume-type functionals of the induced metric, and their
first variation at the constant background is
$\propto\int\eta^{\mu\nu}\langle\bar E_\mu,\partial_\nu\delta\Theta\rangle_\sharp
=-\int\langle\partial_\nu\bar E^\nu,\delta\Theta\rangle_\sharp=0$;
the Einstein term satisfies $R(\mathring\omega(\bar e))=0$ and
$\delta\!\int\!\epsilon\,e\,e\,R
 =\int\!\epsilon\,e\,e\,\mathrm d(\delta\mathring\omega)
 =\int\!\mathrm d(\epsilon\,e\,e\,\delta\mathring\omega)=0$
since $\mathrm d e=0$ for the constant background tetrad.
Both integral statements are verified exactly (SymPy, polynomial
boundary-vanishing variations with linear and quadratic coefficients,
linearised Levi--Civita connection) in
\texttt{tools/verify\_canonical\_spin\_current.py}, check C7.
\end{proof}

\begin{corollary}[Scheme inequivalence at the flat point]
The effective Palatini scheme excludes every flat affine representer
(\texttt{GAP-10T-FLAT-NOGO}), while the gradient-composite scheme admits it
for all coefficients.  The two variational schemes of the same minimal
action are dynamically inequivalent, and the surviving minimal continuation
of canonical UBT is the composite one.
\end{corollary}

\begin{tcolorbox}[colback=orange!6!white,colframe=orange!70!black,title=Scope
and non-claims]
Proved for the gradient-composite torsion-free scheme with $V\equiv0$ and
Lorentz-real variations.  Not proved: the self-consistent $D$-composite
variation, stationarity on curved backgrounds, uniqueness or stability of
the flat point, any derivation of $S_{\rm HP}$ or its coefficients.
\texttt{GAP-10T-DYN} and \texttt{GAP-10D} remain NARROWED.  Note the
$\Lambda$-independence of flat stationarity: the composite scheme imposes
only divergence-type conditions on covariantly constant metric-sector
mismatches, a unimodular-like feature that must be treated explicitly in
the curved dynamical analysis.
\end{tcolorbox}

\end{document}
